% !TeX program = xelatex
% 2026-09-22 批注修订版：以用户本轮提供的文件为基准，仅新增或改写文字标红。
% 标红基准保存于 tmp/revision_20260922/before/；沿用文字与公式恢复黑色。
% 研究内容统一为概述与三项子内容；突出机理、方法与验证之间的对应关系。
% 研究目标及两项科学问题均采用单段论述；反向建模术语全文统一。
% 研究方案保留三组核心公式；可行性分析参照指定面上及重大项目材料重写。
% 可编译主文件；参考文献在本文件内，技术路线图采用正文内的TikZ矢量图。
% 正文按申请书模板的一、二、三部分组织。题目另填模板“项目名称”。
% 申请金额按重点课题上限30万元测算；技术指标为拟定目标。
% 已有工作依据：26-0912.pdf、已有组会材料及已核查的控制程序。
% 匿名研究稿不列作已发表论文；机器人代码仅作只读核查。
% 本项目不采用结构模态建模；文献、已有成果与奖励中的相关名称保留原有表述。
% 地面实验依托现有基座振动模拟与机器人作业实验平台，以基座运动模拟柔性桁架振动。
\documentclass[UTF8,a4paper,zihao=-4,fontset=fandol]{ctexart}
\usepackage[left=2.6cm,right=2.6cm,top=2.5cm,bottom=2.5cm]{geometry}
\usepackage{amsmath,amssymb,bm,booktabs,longtable,array,tabularx}
\usepackage{enumitem,setspace,xcolor,caption,needspace,placeins}
\usepackage{graphicx}
\usepackage{tikz}
\usetikzlibrary{arrows.meta,positioning}
\graphicspath{{./}{../}}
\usepackage[unicode,hidelinks]{hyperref}
\usepackage{xurl}
\hypersetup{pdftitle={柔性桁架支撑下空间机器人的运动控制与接触引导组装研究},pdfsubject={南航重点实验室开放课题申请书报告正文}}
\setstretch{1.32}
\setlength{\parindent}{2em}
\setlength{\parskip}{0.15em}
\setlength{\emergencystretch}{2em}
\clubpenalty=10000
\widowpenalty=10000
\displaywidowpenalty=10000
\raggedbottom
\setlist[enumerate]{leftmargin=2.2em,itemsep=0.3em,topsep=0.3em}
\captionsetup{font=small,labelfont=bf,labelsep=quad,skip=8pt}
\ctexset{
 section={format=\zihao{4}\bfseries,name={,、},number=\chinese{section},aftername={},beforeskip=1.2em,afterskip=0.7em},
 subsection={format=\zihao{-4}\bfseries,number=\arabic{subsection},aftername={.\ },beforeskip=1em,afterskip=0.5em},
 subsubsection={format=\zihao{-4}\bfseries,number=\arabic{subsection}.\arabic{subsubsection},aftername={\quad},beforeskip=0.8em,afterskip=0.3em}
}
\newcommand{\topic}[1]{\par\Needspace{4\baselineskip}\addvspace{0.5em}\noindent\textbf{#1}\par\nobreak}
\AddToHook{cmd/section/before}{\Needspace{9\baselineskip}}
\AddToHook{cmd/subsection/before}{\Needspace{7\baselineskip}}
\AddToHook{cmd/subsubsection/before}{\Needspace{5\baselineskip}}
\newcommand{\needinfo}[1]{\par\noindent{\color{red!65!black}\textbf{〔待补充〕}#1}\par}
\newcommand{\rev}[1]{{\color{red}#1}}
\pdfstringdefDisableCommands{\def\rev#1{#1}}
\newcommand{\doi}[1]{\href{https://doi.org/#1}{\nolinkurl{#1}}}
\newcommand{\SE}{\mathrm{SE}(3)}
\newcolumntype{Y}{>{\raggedright\arraybackslash}X}
\newcolumntype{P}[1]{>{\raggedright\arraybackslash}p{#1}}
\renewcommand{\arraystretch}{1.25}
\pagestyle{plain}

\begin{document}

\begin{center}
 {\zihao{3}\bfseries 柔性桁架支撑下空间机器人的运动控制与接触引导组装研究\par
\par}
 \vspace{0.55em}
 {\zihao{-4}航空航天结构力学及控制全国重点实验室开放课题\quad 重点课题}
\end{center}
\vspace{0.35em}

% ==================== 报告正文开始 ====================
\section{立项依据与研究内容}

\subsection{项目的立项依据}

\subsubsection{研究意义}

大型空间天线、太空望远镜等装备的发展，对空间设施的尺度、功能与建造精度提出了更高要求。受运载包络和发射能力限制，依靠单次整体发射与在轨展开难以满足建造需求，通过模块分批运输\rev{和}机器人在轨转运\rev{、}装配逐步形成大型设施，成为拓展空间装备能力的重要途径\cite{hu2025,wu2024}。在这一过程中，机器人需要依附已建结构完成模块拾取、转运、接口对准\rev{与}安装\rev{，基座}运动、\rev{负载}变化\rev{和}接触状态转换使\rev{其}面临复杂的动力学条件，\rev{亟需发展}兼顾运动精度与接触柔顺性的控制方法。

机器人依附柔性桁架作业时，机械臂运动产生的反作用力与力矩经基座\rev{激发}桁架振动，桁架振动又引起基座\rev{位姿}变化并传递至作业端，\rev{形成双向动力学耦合}；惯性参数、负载及执行时延等偏差进一步\rev{影响}控制\rev{精度}。提高反馈增益\rev{虽可}减小\rev{部分}运动误差，\rev{却}限制了机器人顺应接触约束的调整能力，\rev{接口}位姿偏差与\rev{基座}振动易引起接触\rev{冲击或}卡滞。\rev{同时}，高维\rev{耦合}模型\rev{的}在线求解难以兼顾描述精度与反馈实时性，需要提取影响末端运动\rev{和受力}的主要\rev{界面}信息，\rev{分析驱动约束下实际阻抗}的\rev{形成规律}，并\rev{利用}接触\rev{几何修正对准目标}。

本项目以柔性桁架支撑的空间机器人为研究对象，研究\rev{安装界面耦合作用}的\rev{反向}动力学表征，结合几何阻抗与动力学优化开展\rev{受约束}运动控制，利用接触\rev{场配准}和\rev{有界参考修正}实现柔顺组装，并通过地面\rev{实验}验证。研究\rev{将揭示界面扰动、目标估计}与\rev{受约束执行对末端响应的}共同\rev{影响}，为大型空间\rev{设施}机器人在轨组装提供理论依据与\rev{方法}支撑。

\subsubsection{国内外研究现状及发展动态}

\topic{（1）柔性结构与空间机器人的耦合动力学}

\rev{在轨}操作由航天器捕获向大型结构组装拓展，\rev{推动}动力学研究由基座与机械臂的运动耦合，转向机器人运动、结构变形与组装过程的共同作用\cite{hu2025,wu2024,papadopoulos2021}。机器人所依附\rev{结构}的质量分布、连接关系和边界随组装进程变化，模型需要\rev{兼顾}构型演化\rev{与}作业状态转换。

现有方法主要\rev{在}动态特征保留与模型规模\rev{之间寻求平衡}。模态缩减\rev{和}模型降阶通过保留主要变形特征降低系统维数\cite{sonneville2021,min2026}；时变边界研究\rev{进一步}表明，模型参数与约束\rev{需}随作业状态调整\cite{ni2024}。超大型结构的组装分析\rev{揭示了}科里奥利效应等空间扰动\rev{的}影响\cite{yang2023}，柔性模块接触\rev{模型}则将冲击、结构变形与机器人运动联系起来\cite{wucontact2026}。\rev{这些研究说明}，\rev{面向}组装\rev{控制的}模型\rev{应同时关注安装端}运动\rev{传递、}接触约束建立\rev{和}载荷\rev{变化。}

模型应用已由\rev{响应}预测\rev{延伸到}运动与控制\rev{设计}。振动反馈\rev{利用}机器人\rev{与}结构\rev{的}相互作用调节响应\cite{swei2020}；基于线性分式表示的路径优化将柔性动态\rev{和}参数不确定性纳入性能评价\cite{rodrigues2025}；直接轨迹优化\rev{与}递推动力学方法\rev{进一步结合}耦合方程\rev{和}作业约束\rev{改善}求解效率\cite{liu2026,gong2026}。但高保真仿真、轨迹规划与高频反馈对模型信息的要求\rev{不同}，计算效率的提高还需结合测量\rev{误差}评价其闭环\rev{效果}。面向实时控制，\rev{仍需}提取安装\rev{界面}运动与受力对\rev{末端}响应的主要影响，\rev{明确简化模型}的\rev{误差及适用范围}。

\topic{（2）运动振动协调与阻抗柔顺控制}

运动振动协调\rev{要求}在完成末端任务的同时控制机器人对柔性结构的反作用。反作用零空间及其优化方法利用冗余运动减少结构激励\cite{nenchev1999,yin2023}\rev{，}柔性状态估计与振动反馈则\rev{依据}结构响应进行补偿或主动调节\cite{patel2024,zhou2025}。前者受构型\rev{和}剩余自由度限制，后者依赖状态信息\rev{与}可用执行机构，需要结合具体\rev{支撑与驱动}条件\rev{设计}。

接触任务进一步\rev{要求}位置、姿态\rev{及}力和力矩具有一致的几何表达。基于$\SE$的几何阻抗\rev{通过}刚体位姿空间中\rev{的}势能\rev{与误差描述协调平移}和\rev{转动}\cite{seo2023}；李群与李代数势能比较\rev{及}阻抗映射研究，\rev{揭示了}误差表示\rev{对}有限位姿偏差下恢复作用\rev{的影响}\cite{seo2024,kim2025}。\rev{针对模型与环境不确定性，滑模补偿和模糊增益调节用于减弱动力学失配}\cite{liudb2024}\rev{，环境刚度估计与分层控制进一步结合柔性关节、基座振动及接触变化协调响应}\cite{liuadaptive2026}\rev{。}

\rev{上述研究}为\rev{振动环境下的柔顺}控制提供了基础，但设计阻抗仍需通过\rev{有限}关节驱动实现\rev{。关节运动或力矩约束激活时}，运动跟踪、反作用调节与柔顺目标可能\rev{相互制约}；基座运动又改变接触端的实际速度与受力。因此，\rev{需将几何阻抗、动力学分配与安装端作用统一分析，}研究驱动约束下\rev{阻抗响应的}可实现性。

\topic{（3）接触信息引导的位姿估计与柔顺组装}

柔顺组装正由利用接触反馈调节受力，\rev{向}利用接触约束修正目标\rev{位姿发展}。已有空间模块对接研究\rev{通过}接触力误差调整柔顺轨迹，并\rev{完成}地面验证\cite{shi2024}。\rev{但}较小\rev{的}接触力\rev{仍可能对应}错误的接口偏置，目标位姿不确定\rev{时}，需要同时利用\rev{受力反馈}与配合几何。

接触几何利用\rev{主要包括}联合状态估计\rev{与}接触模型配准。前者\rev{融合}触觉、机器人运动和接触约束，\rev{借助}主动运动更新物体位姿与接触状态\cite{kimtactile2023}；后者通过探测采样局部接触流形，与\rev{预建}接口模型配准获得位姿修正\cite{negi2025}。柔顺运动与主动接触交互的结合，使机器人能够逐步定位孔位\rev{和}修正插接方向\cite{chen2025}。\rev{这些}方法\rev{通过多次观测累积几何}信息，\rev{其估计能力取决于接口形状、接触分布}与\rev{探测运动}。

\rev{在}柔性支撑条件下，接触位姿变化同时\rev{包含}目标偏差、基座振动\rev{及}柔顺响应\rev{，}用于估计的探测运动和参考更新又会改变接触\rev{载荷}。因此，\rev{需分析基座测量与时序误差向目标配准}的\rev{传播}，协调参考修正与实际\rev{跟踪能力，明确估计精度改善能够转化为对准精度和接触载荷改善的条件}。

综上，\rev{现有}研究已形成耦合动力学、运动振动协调\rev{及}接触估计的方法基础。\rev{本项目进一步贯通安装界面表征、受约束阻抗实现与接触引导，研究各环节误差对组装响应的共同影响，并通过可控基座激励}下的地面\rev{实验检验其适用范围}。

\subsubsection{与相关指南的关系}

本项目主要对应指南2.2“轻质柔性结构动力学设计与分析”中柔性空间结构在轨组装及刚柔耦合体系力学控制研究，围绕机器人安装界面的动态作用、受约束运动与接触组装开展研究；同时衔接指南2.5关于多源不确定性条件下闭环分析与控制的研究内容。
\begingroup
\small
\setstretch{1.08}
\setlength{\parskip}{0pt}

\renewcommand{\refname}{主要参考文献}
\begin{thebibliography}{99}
\setlength{\itemsep}{0.15em}
\bibitem{hu2025} 胡海岩, 田强, 文浩, 罗凯, 马小飞. 极大空间结构在轨组装的动力学与控制[J]. 力学进展, 2025, 55(1): 1--29. DOI: \doi{10.6052/1000-0992-24-044}.
\bibitem{wu2024} 吴志刚, 蒋建平, 邬树楠, 李庆军, 王兴, 谭述君, 邓子辰. 航天结构空间组装动力学与控制研究进展[J]. 力学进展, 2024, 54(2): 344--390. DOI: \doi{10.6052/1000-0992-23-041}.
\bibitem{papadopoulos2021} Papadopoulos E, Aghili F, Ma O, Lampariello R. Robotic manipulation and capture in space: A survey[J]. Frontiers in Robotics and AI, 2021, 8: 686723. DOI: \doi{10.3389/frobt.2021.686723}.
\bibitem{sonneville2021} Sonneville V, Scapolan M, Shan M, Bauchau O A. Modal reduction procedures for flexible multibody dynamics[J]. Multibody System Dynamics, 2021, 51(4): 377--418. DOI: \doi{10.1007/s11044-020-09770-w}.
\bibitem{min2026} Min K, Fan W, Ren H, Zhou P, Wang Y, Zhang X. Dynamic modeling and model order-reduction of on-orbit assembled multi-plate structures[J]. International Journal of Mechanical Sciences, 2026, 313: 111296. DOI: \doi{10.1016/j.ijmecsci.2026.111296}.
\bibitem{ni2024} Ni S, Chen W, Chen T. Dynamic modeling and analysis of multi-flexible-link space manipulators under time-varying dynamic boundary conditions[J]. Advances in Space Research, 2024, 74(10): 5224--5243. DOI: \doi{10.1016/j.asr.2024.07.067}.
\bibitem{yang2023} Yang G, Zhang L, Yu S, Meng S, Wang Q, Li Q. Influences of space perturbations on robotic assembly process of ultra-large structures[J]. Nonlinear Dynamics, 2023, 111(11): 10025--10048. DOI: \doi{10.1007/s11071-023-08395-w}.
\bibitem{wucontact2026} Wu L, Wang M, Luo J, Zhang D. Dynamic modeling and analysis of space robots with an irrotational contact model for on-orbit assembly[J]. Aerospace Science and Technology, 2026, 176: 112117. DOI: \doi{10.1016/j.ast.2026.112117}.
\bibitem{swei2020} Swei S S M, Jenett B, Cramer N B, Cheung K. Modeling and control of robot-structure coupling during in-space structure assembly[C]//AIAA Scitech 2020 Forum. 2020: AIAA 2020-1544. DOI: \doi{10.2514/6.2020-1544}.
\bibitem{rodrigues2025} Rodrigues R, Preda V, Sanfedino F, Alazard D. Dynamics modeling and path optimization for the on-orbit assembly of large flexible structures using a multi-arm robot[J/OL]. CEAS Space Journal, 2025. DOI: \doi{10.1007/s12567-025-00675-y}.
\bibitem{liu2026} Liu P, Wang M, Ma W, Luo J, Liu C. Direct trajectory optimization of coupled constrained system formed by the space manipulator and the flexible structure[J]. Aerospace Science and Technology, 2026, 168: 111038. DOI: \doi{10.1016/j.ast.2025.111038}.
\bibitem{gong2026} Gong Y, Wen H. Recursive modeling and trajectory optimization of space manipulator mounted on flexible structure[J]. Applied Mathematical Modelling, 2026, 155: 116741. DOI: \doi{10.1016/j.apm.2026.116741}.
\bibitem{nenchev1999} Nenchev D N, Yoshida K, Vichitkulsawat P, Uchiyama M. Reaction null-space control of flexible structure mounted manipulator systems[J]. IEEE Transactions on Robotics and Automation, 1999, 15(6): 1011--1023. DOI: \doi{10.1109/70.817666}.
\bibitem{yin2023} 尹旺, 王翔, 王为, 刘冬雨. 柔性基座冗余机械臂的最优反作用控制[J]. 中国空间科学技术, 2023, 43(2): 144--154. DOI: \doi{10.16708/j.cnki.1000-758X.2023.0029}.
\bibitem{patel2024} Patel D, Damaren C J. Tip and vibration control of space robots using estimated flexible coordinates[J]. Frontiers in Space Technologies, 2024, 5: 1447545. DOI: \doi{10.3389/frspt.2024.1447545}.
\bibitem{zhou2025} Zhou W, Ye Z, Wu S, Li Y. Coupling dynamic modeling and vibration control of quadruped-robot and large space structure during on-orbit assembly[J]. Space Solar Power and Wireless Transmission, 2025, 2(1): 1--9. DOI: \doi{10.1016/j.sspwt.2025.02.002}.
\bibitem{seo2023} Seo J, Prakash N P S, Rose A, Choi J, Horowitz R. Geometric impedance control on $\SE$ for robotic manipulators[J]. IFAC-PapersOnLine, 2023, 56(2): 276--283. DOI: \doi{10.1016/j.ifacol.2023.10.1581}.
\bibitem{seo2024} Seo J, Prakash N P S, Choi J, Horowitz R. A comparison between Lie group- and Lie algebra-based potential functions for geometric impedance control[C]//2024 American Control Conference (ACC). 2024: 1335--1342. DOI: \doi{10.23919/ACC60939.2024.10644201}.
\bibitem{kim2025} Kim J, Sung M, Choi Y, Park J, Chung W K. Impedance control design framework using commutative map between $\SE$ and $\mathfrak{se}(3)$[J]. IEEE Transactions on Robotics, 2025, 41: 6193--6212. DOI: \doi{10.1109/TRO.2025.3619066}.
\bibitem{liudb2024} 刘东博, 陈力. 面向在轨组装服务的空间机器人动力学分析及抑振阻抗控制设计[J]. 力学季刊, 2024, 45(3): 688--696. DOI: \doi{10.15959/j.cnki.0254-0053.2024.03.008}.
\bibitem{liuadaptive2026} Liu P, Ma W, Wang M, Luo J, Liu C. Adaptive impedance control of a flexible-joint manipulator for space assembly operations under base vibrations and environmental uncertainty[J]. Advances in Space Research, 2026, 77(7): 8057--8074. DOI: \doi{10.1016/j.asr.2026.02.005}.
\bibitem{shi2024} 史玲玲, 肖晓龙, 张晓峰, 范立佳, 单明贺, 田强. 空间机器人在轨组装多模块单元的对接力控制与地面实验[J]. 力学学报, 2024, 56(3): 800--816. DOI: \doi{10.6052/0459-1879-23-458}.
\bibitem{kimtactile2023} Kim S, Jha D K, Romeres D, Patre P, Rodriguez A. Simultaneous tactile estimation and control of extrinsic contact[C]//2023 IEEE International Conference on Robotics and Automation (ICRA). 2023: 12563--12569. DOI: \doi{10.1109/ICRA48891.2023.10161158}.
\bibitem{negi2025} Negi A, Manyar O M, Penmetsa D K, Gupta S K. Accurate pose estimation using contact manifold sampling for safe peg-in-hole insertion of complex geometries[C]//2025 IEEE 21st International Conference on Automation Science and Engineering (CASE). 2025: 617--624. DOI: \doi{10.1109/CASE58245.2025.11163996}.
\bibitem{chen2025} Chen Y, Kimble K, Qian H H, Chanrungmaneekul P, Seney R, Hang K. Robust peg-in-hole assembly under uncertainties via compliant and interactive contact-rich manipulation[C]//Robotics: Science and Systems XXI. 2025. DOI: \doi{10.15607/RSS.2025.XXI.060}.
\bibitem{qi2025} 祁伊凡, 单明贺, 田强. 空间多柔体系统缩比等效动力学建模方法研究[J]. 中国科学: 物理学 力学 天文学, 2025, 55(2): 224518. DOI: \doi{10.1360/SSPMA-2024-0302}.
\end{thebibliography}
\endgroup

\subsection{项目的研究内容、研究目标及拟解决的关键科学问题}

\subsubsection{研究内容}

本项目面向柔性\rev{桁架}支撑下机器人的精确运动与模块组装，采用理论分析、数值仿真与地面\rev{实验}相结合的方法，\rev{围绕安装界面动力学传递、受约束阻抗响应及接触估计与执行协调，开展以下三项}研究。

\Needspace{8\baselineskip}
\topic{（1）柔性桁架支撑下机械臂的耦合动力学建模与运动控制}

\rev{针对}桁架\rev{振动与}机械臂运动的双向耦合\rev{，以及关节驱动约束引起}的\rev{阻抗实现偏差}，从安装界面的运动与受力出发，\rev{研究耦合作用}的\rev{反向}动力学\rev{表征与}几何阻抗\rev{优化控制，揭示运动精度、接触柔顺性}与\rev{安装端反作用之间}的\rev{关联}。主要研究内容如下：
\begin{enumerate}[label=\arabic*)]
 \item 研究安装\rev{界面}运动与受力\rev{的功共轭}表征，建立以作业端为浮动基的\rev{反向}动力学模型，\rev{揭示}基座振动\rev{与关节}运动对末端响应的\rev{耦合传递}规律。
 \item 研究\rev{$\SE$}几何阻抗与动力学优化\rev{的协同设计，将关节运动及力矩约束纳入统一求解，形成兼顾位姿跟踪与柔顺响应的受约束运动}控制方法。
 \item 研究\rev{界面载荷}建模偏差\rev{与}约束\rev{激活下的}阻抗\rev{可实现性}，\rev{分析局部闭环稳定与误差有界条件，阐明}运动精度、柔顺\rev{响应及}安装端反作用的关联。
\end{enumerate}

\Needspace{8\baselineskip}
\topic{（2）目标位姿不确定条件下的接触对准与柔顺组装}

\rev{针对}基座运动\rev{引起的}接触观测\rev{偏差}与参考\rev{执行失配，以接触流形的几何约束}和接触\rev{反馈}为\rev{依据}，研究\rev{目标位姿的局部可辨识性}与\rev{有界}参考修正\rev{，揭示接触估计、参考更新和实际阻抗共同作用下}的\rev{对准}与\rev{受力规律}。主要研究内容如下：
\begin{enumerate}[label=\arabic*)]
 \item \rev{研究}接触\rev{流形约束下六维目标位姿的局部可辨识性}，分析基座测量\rev{与时序}误差\rev{向}接触\rev{场}配准的\rev{传播规律}，\rev{建立弱}约束\rev{方向}的\rev{先验融合方法}。
 \item \rev{研究}五轴约束导纳\rev{与轴向接近运动的协调机制}，\rev{保持参考端口耗散性与}零输入保持性质，\rev{构建}修正速率\rev{和累计}偏移\rev{受限}的\rev{柔顺组装策略}。
 \item \rev{研究}目标估计\rev{、参考修正与受约束执行的耦合误差传递，阐明估计}精度\rev{转化为对准精度和}接触载荷\rev{改善的条件}，\rev{形成}引导\rev{与执行协调方法}。
\end{enumerate}

\Needspace{8\baselineskip}
\topic{（3）柔性支撑条件下机器人运动与组装的地面实验验证}

\rev{针对振动环境下运动精度与接触响应的实验验证需求，}依托基座振动模拟与机器人作业实验平台，\rev{研究安装界面}运动\rev{复现与多源}同步\rev{观测方法}，通过\rev{分项对照和组装集成实验}，\rev{检验耦合}动力学描述\rev{、受约束}控制\rev{及接触引导}方法\rev{的有效性}。主要研究内容如下：
\begin{enumerate}[label=\arabic*)]
 \item \rev{研究安装界面动态边界的运动复现}与\rev{载荷表征，构建多源同步观测}与\rev{独立精度}评价方法\rev{，明确基座激励对桁架振动影响的可表征范围}。
 \item 开展位姿保持与轨迹跟踪\rev{对照实验}，\rev{辨析界面载荷建模偏差}与\rev{约束激活对}运动精度\rev{和}实际阻抗\rev{的影响，验证动力学模型}与\rev{控制的实时性}。
 \item 开展典型接口对准、插接\rev{与}到位保持\rev{集成实验}，\rev{分离}接触配准\rev{和}参考修正的\rev{性能贡献}，评价振动环境下\rev{组装精度、接触载荷}与\rev{方法适用范围}。
\end{enumerate}

\Needspace{5\baselineskip}\subsubsection{研究目标}

面向大型空间结构在轨组装需求，建立基于安装界面运动与受力\rev{表征}的\rev{反向}动力学\rev{模型}，揭示桁架振动\rev{与机械臂运动的耦合传递规律及}驱动约束\rev{下阻抗}响应的\rev{形成机制}，形成几何阻抗与动力学优化相结合的运动控制方法；\rev{阐明接触几何约束下目标位姿的}局部\rev{可辨识性及估计与执行}误差对组装响应的\rev{共同}影响，\rev{形成}接触\rev{场配准、有界}参考修正与\rev{受约束运动控制相协调}的柔顺组装方法；\rev{建立安装端运动复现}与\rev{独立测量评价的地面}实验\rev{方法}，\rev{完成运动控制和典型接口组装}验证\rev{，明确}所提方法在不同基座振动、末端负载及初始位姿偏差条件下的适用范围。

\subsubsection{拟解决的关键科学问题}

\topic{（1）机器人与柔性桁架耦合作用下受约束运动与阻抗响应的形成机制}

柔性桁架振动通过安装端运动与受力改变机械臂的动态响应，机械臂驱动又通过安装端反作用影响柔性桁架的振动响应，使作业端运动与安装端受力相互关联。面向实时控制，需要保留影响末端响应的主要耦合信息，而安装端作用力与力矩的近似描述会带来动力学偏差；关节运动范围与驱动力矩限制还可能使期望运动无法完全实现，导致实际阻抗响应偏离设计值。因此，安装端耦合作用如何经受约束力矩分配影响末端运动与柔顺响应，是本项目需要解决的关键科学问题。


\topic{（2）基座运动\rev{与接触约束共同作用}下\rev{对准精度与柔顺组装响应的形成机制}}

\rev{基座运动通过观测坐标与测量时序影响}接触\rev{位姿}估计\rev{，接触几何与观测分布又决定不同位姿方向的约束能力，使目标估计}与参考修正\rev{相互关联。面向目标位姿不确定的}组装\rev{过程，需要利用接触信息持续修正对准目标，而参考更新受到关节驱动约束与实际阻抗}响应的\rev{共同限制；}估计\rev{偏差}和跟踪\rev{滞后还可能转化为额外接触载荷}，\rev{使目标}估计精度的改善\rev{难以}直接\rev{对应于组装性能提升。因此，接触估计与参考修正如何经受约束运动执行影响对准精度与接触柔顺性，是本项目需要解决}的\rev{关键科学问题}。

\subsection{拟采取的研究方案及可行性分析}

\subsubsection{总体研究思路与技术路线}

对应三项研究内容，按照“\rev{界面}耦合\rev{表征}与运动控制、接触引导\rev{与受约束执行协调}、地面\rev{实验}验证”的\rev{技术路线展开}。首先\rev{，}以安装端\rev{运动}与\rev{载荷建立反向动力学描述}，结合几何阻抗与动力学优化\rev{研究驱动约束下的末端响应；其次}，利用接触流形的\rev{局部}几何\rev{信息估计目标位姿}，\rev{分析观测误差传播并}通过有界参考修正\rev{协调}对准与\rev{接触受力}；最后\rev{，复现安装端运动}，开展\rev{模型校核、}分项对照与\rev{组装}集成\rev{实验，以实测结果检验误差分析和参数选取依据}。技术路线如图\ref{fig:route}所示。

\begin{figure}[htbp]
\centering
\begin{tikzpicture}[
 block/.style={draw=blue!45!black,fill=blue!3,rounded corners=3pt,
   text width=12.4cm,align=center,inner sep=8pt,font=\small},
 flow/.style={-{Stealth[length=2.5mm]},thick,draw=blue!50!black},
 node distance=0.55cm]
 \node[block,fill=blue!8] (need) {\textbf{柔性桁架上的精确运动与柔顺组装}\\[3pt]
 基座振动\qquad 关节驱动约束\qquad 目标位姿偏差};
 \node[block,below=of need] (control) {\textbf{研究内容一\quad 耦合动力学与运动控制}\\[3pt]
 \rev{反向}动力学\quad 几何阻抗\quad 动力学优化};
 \node[block,below=of control,draw=green!40!black,fill=green!3] (contact) {\textbf{研究内容二\quad 接触对准与柔顺组装}\\[3pt]
 接触流形表征\quad 接触场配准\quad 有界参考修正};
 \node[block,below=of contact,draw=violet!60!black,fill=violet!3] (test) {\textbf{研究内容三\quad 地面\rev{实验}验证}\\[3pt]
 \rev{安装端运动复现}与同步\rev{观测}\\[2pt]
 位姿保持\quad 轨迹跟踪\quad 典型接口组装};
 \draw[flow] (need) -- (control);
 \draw[flow] (control) -- (contact);
 \draw[flow] (contact) -- (test);
\end{tikzpicture}
\caption{项目技术路线与验证关系}
\label{fig:route}
\end{figure}
\FloatBarrier

\subsubsection{\rev{反向}动力学建模与运动控制}

\rev{以}作业端为浮动基\rev{、}实际安装\rev{端为}虚拟末端，\rev{沿作业端向安装端重新组织运动链，建立反向动力学模型。}通过\rev{安装界面}的\rev{运动}与\rev{受力计入}桁架对机械臂的作用，\rev{分析}基座\rev{振动}、关节运动与\rev{末端负载的耦合传递。模型基本形式为：}

{\color{red}\begin{equation}
 \bm M\dot{\bm\nu}+\bm h
 =\bm S^{\mathsf T}\bm\tau_j
 +\bm J_b^{\mathsf T}\bm F_b
 +\bm J_g^{\mathsf T}\bm F_e .
 \label{eq:reverse_dynamics}
\end{equation}}

\rev{其中，$\bm\nu$由作业端体速度和关节速度组成，}$\bm M$\rev{与}$\bm h$\rev{分别为惯性矩阵和包含科氏、离心及重力作用的偏置项；$\bm S$为关节选择矩阵，}$\bm\tau_j$\rev{为关节驱动}力矩\rev{；}$\bm F_b$\rev{、}$\bm F_e$\rev{及其雅可比矩阵$\bm J_b$、}$\bm J_g$\rev{分别对应实际作业端和}安装端\rev{。自由运动时}$\bm F_b=0$\rev{，接触时由力传感测量计入；力旋量}与\rev{速度采用功共轭坐标表达。}

以现有\rev{安装端}静力平衡近似为基线，结合\rev{实测基座}运动与\rev{界面载荷}，研究动态\rev{偏差随激励和负载}的\rev{变化规律，形成}控制\rev{所需}的\rev{载荷描述}与\rev{误差范围}。地面实验中，\rev{模拟机械臂末端与}作业机械臂基座固定连接，对应模型的实际安装端；\rev{通过指令、实测}运动与\rev{受力对照}，校核地面激励\rev{和}控制模型的对应关系。

\begingroup\tolerance=9999\emergencystretch=8em
\rev{在}$\SE$\rev{上统一描述}实际\rev{位姿}$T_b$与期望位姿$T_d$，采用\rev{局部}误差\rev{坐标}\allowbreak$\color{red}\bm\lambda=\operatorname{Log}(T_b^{-1}T_d)^\vee$，\rev{将}几何阻抗\rev{映射}为误差动态\cite{kim2025}：
\par\endgroup

\begin{equation}
 \bm A_\lambda\ddot{\bm\lambda}
 +\bm D_\lambda\dot{\bm\lambda}
 +\bm K_\lambda\bm\lambda
 =\bm\gamma .
 \label{eq:impedance}
\end{equation}

\rev{其中}，$\bm A_\lambda$、$\bm D_\lambda$\rev{、}$\bm K_\lambda$为\rev{经几何映射得到的等效}惯量、阻尼和刚度矩阵，$\bm\gamma$为\rev{与误差坐标及}相对运动方向\rev{一致}的\rev{广义接触输入}。\rev{结合轨迹前馈和}速度映射\rev{生成作业端}参考加速度，校核\rev{误差坐标与其时间导数}的\rev{一致性}。

将\rev{参考加速度嵌入}动力学二次规划\rev{（QP）}，以\rev{参考跟踪、}关节虚拟阻尼和力矩正则为目标，\rev{以反向}动力学\rev{为}等式约束，\rev{统一计入}关节\rev{运动范围、力矩及其变化率限制}。\rev{通过}参考\rev{松弛处理期望运动}不可实现\rev{的情形}，\rev{研究约束激活、优化权重与}实际阻抗的\rev{关系}。\rev{以几何阻抗误差能量为基础，}在\rev{优化}可行\rev{、测量}误差有界\rev{及参考足够平滑}的条件下，分析\rev{典型约束激活区间内}的局部稳定与误差界，结合\rev{受控扰动实验确定阻抗参数和参考变化速率}的\rev{选取依据}。

\subsubsection{基于接触流形表征与参考修正的柔顺组装}

\rev{以}已有物理方向辅助接触场配准方法（PCF）\rev{为基础}，以离线接触位姿约束场值，以功共轭变换\rev{后}的物理方向约束局部斜率，结合全局场与局部拟合表征接口接触\rev{流形}。在线以目标先验为初值，\rev{在统一参考系中}求解六维\rev{位姿}：
\begin{equation}
 \widehat T_\star=\underset{T\in\SE}{\arg\min}
 \sum_i\rho\!\left(F(T^{-1}T_i)\right).
 \label{eq:registration}
\end{equation}
其中，$T_i$为接触位姿，$F$为局部位姿坐标中的接触高度残差，$\rho$为鲁棒损失。\rev{结合}实测基座位姿\rev{和时序}修正观测坐标，\rev{按接口特征尺度统一平移与转动的}局部\rev{度量，利用残差雅可比的奇异值及对应方向分析局部可辨识性，研究弱约束方向的先验融合}与执行限幅\rev{，并}通过扰动回放\rev{检验观测}误差\rev{传播}。

\rev{利用}约束导纳捕获方法（RAC）调节两个横向\rev{与}三个姿态通道，轴向接近单独处理。在正定\rev{导纳}阻尼\rev{和}包含零的闭凸增量\rev{约束下}，\rev{依据共轭接触输入求解修正增量，限制变化速率与}累计偏移，\rev{并校核}有限转动步长，\rev{保持}零输入保持\rev{性质及}参考端口离散广义功\rev{关系}。\rev{根据估计敏感程度设置导纳阻尼和}修正\rev{边界，将}位姿\rev{及其}速度、加速度参考\rev{一致地传递至下层控制。}

\rev{通过局部线性化与受控扰动}，分析\rev{目标估计、参考}修正\rev{及实际}跟踪误差对到位精度和接触\rev{载荷的贡献，考察加速度松弛对误差传递}的影响，\rev{形成}参考修正与\rev{驱动}能力的匹配\rev{方法}。局部\rev{分析}和参考子系统\rev{功关系用于解释作用机制}，整机\rev{性能通过配准、参考修正及联合控制的}对照\rev{实验}评价。

\subsubsection{地面\rev{实验}系统与分项、集成验证}

依托现有平台，\rev{由}基座运动模拟\rev{机械臂}驱动作业机械臂基座产生可重复的平移\rev{与}转动，\rev{复现}桁架振动引起的安装端运动。\rev{按}平台工作范围和预实验结果\rev{设置}激励幅值、频率与方向，\rev{以}末端配重调整\rev{负载，}以同一安装状态\rev{下的静止基座}作为基线。同步采集\rev{基座与末端位姿、}关节状态\rev{及两端}力和力矩，\rev{分析激励传递的幅值和相位偏差，建立界面运动与载荷表征方法。传感器}标定\rev{、}外参\rev{测定}和\rev{时序同步作为测量保障}，\rev{末端精度}采用独立于控制反馈的测量评价。

\rev{运动控制验证采用}位姿保持\rev{、}圆轨迹跟踪\rev{和接触扰动实验}，\rev{结合}仿真与实测\rev{对照}，\rev{分析载荷描述}及约束激活\rev{对实际阻抗}的\rev{影响}。\rev{保持任务参考和激励一致，核查实测基座运动，}记录位置均方根误差、旋转测地误差、关节力矩\rev{及其}变化率\rev{、}安装端\rev{受力，并比较设计阻抗与}力和\rev{运动实测响应}；\rev{统计}计算耗时分布、最大值\rev{和超}周期比例，验证实时执行能力。\rev{每组正式对照}至少\rev{重复}5次\rev{，}参数对照\rev{集中}在代表性工况。

\rev{组装验证选取}一种典型接口和代表性载荷，将目标固定于独立刚性基准。先在基座静止时\rev{设置}基础柔顺控制、仅配准、仅参考修正\rev{和}二者结合的四组对照，再\rev{在代表性}平移或转动激励\rev{下}重复关键对照。统一初始偏差、执行\rev{时间}和保持窗口，记录\rev{原始}估计\rev{、限幅参考和实际到位结果，分离估计与执行}误差\rev{，评价}到位\rev{精度}、达标次数、接触力峰值及安装端\rev{反作用}。

\rev{集成验证贯通接近、对准、插接}与\rev{到位保持。正式对照前依据}接口容差\rev{和}测量精度\rev{固定工况及判据}，成功与失败\rev{实验}均纳入评价，调参数据不计入独立验证，报告绝对误差\rev{及}重复\rev{实验}离散程度\rev{。给定基座激励用于检验末端对振动扰动的适应能力}；安装端\rev{反作用}用于分析\rev{机器人}对\rev{支撑结构}的\rev{载荷}作用，其\rev{降低}不直接等同于\rev{桁架振幅降低}。

\subsubsection{可行性分析}

\topic{（1）\rev{研究思路与理论}方法\rev{可行}}

\rev{本项目以安装界面耦合作用为基础，将受约束阻抗实现与接触引导衔接，形成理论分析、数值仿真和地面实验相互验证的研究思路，各项方法具有明确依据。}

\textbf{\rev{在耦合动力学建模方面：}}\rev{桁架与机械臂通过安装界面交换运动与载荷，功共轭表征为界面建模提供了力学基础。团队已开展柔性多体系统模型简化与缩比等效建模，并实现反向机械臂模型。因此，以}现有\rev{静力平衡近似为起点，利用基座运动与安装端载荷分析动态偏差，可为实时控制所需的耦合描述提供可实施的途径。}

\textbf{\rev{在受约束运动控制方面：}}\rev{几何阻抗统一描述位置与姿态响应，动力学二次规划将运动目标与驱动边界纳入同一求解。团队已将两者用于}位姿保持\rev{和}圆轨迹跟踪，\rev{具备力矩控制}与\rev{数据回放基础}。在此基础上\rev{分析}建模偏差\rev{、参考松弛和约束激活，通过局部误差动态}及\rev{受控扰动实验检验阻抗可实现性}，\rev{研究思路可行}。

\textbf{\rev{在}接触\rev{引导组装方面：}}\rev{接触流形为目标估计提供几何约束，约束导纳将受力反馈转化为有界修正。}已有PCF配准与RAC捕获完成了仿真、真机配准回放及多轮对接验证，可\rev{支撑局部可辨识性、误差传播}与参考\rev{端口功关系分析。结合基座运动}修正\rev{观测}，并由\rev{静止}基座逐步扩展\rev{到振动工况，可分步检验参考修正与力矩控制}的\rev{协调效果}。

\topic{（\rev{2}）\rev{实验方案与}研究\rev{条件可行}}

\textbf{\rev{从实验方案看：}}模拟\rev{机械臂与}作业机械臂\rev{的固定连接可形成可控、可重复的}安装端\rev{运动环境，配重与接口夹具可设置负载和初始偏差。通过基座与末端位姿、关节状态及两端}受力\rev{的同步观测，可分别评价模型误差、跟踪性能与接触响应；结合独立刚性基准和精度测量，能够对照检验各控制环节，覆盖主要研究问题。}

\textbf{\rev{从研究基础与条件看：}}\rev{团队已有基座振动模拟与机器人作业平台、开放力矩控制程序、运动捕捉系统、IMU和六维力测量条件，并积累了多模块对接实验经验。拟完善的连接组件、测点、时序同步和接口夹具可依托现有设备、共享条件与项目预算落实。按单项模型与控制}校核\rev{、接口对照和振动环境集成验证逐步推进}，\rev{工作量}与两年\rev{研究周期}相\rev{匹配}。

\subsection{本项目的特色与创新之处}

\rev{本项目以安装界面的运动与受力贯通耦合动力学、几何阻抗和接触组装，结合可控基座激励与独立测量，研究柔性支撑条件下运动精度与接触柔顺性的协调关系。特色与创新主要体现在以下两个方面。}

\topic{（1）安装界面\rev{耦合作用与驱动约束共同作用下}的\rev{阻抗实现方法}}

\rev{针对}柔性支撑条件下\rev{设计阻抗与实际响应不一致的问题，在反向动力学框架中将安装界面运动与载荷引入作业端响应分析，揭示界面载荷建模偏差与驱动约束激活对阻抗特性的共同影响；将几何阻抗目标与受约束力矩分配相结合，研究阻抗可实现性及局部误差界，形成依据界面扰动和驱动能力协调运动精度与柔顺响应}的控制方法。

\topic{（2）\rev{基座运动}下接触几何估计与受约束执行\rev{相协调的柔顺组装方法}}

\rev{针对目标估计精度提升难以直接转化为组装性能改善的问题，将接触流形的局部可辨识性与实际阻抗响应}联系起来，\rev{揭示观测误差、}参考修正\rev{和执行偏差向对准误差及接触载荷的传递规律；结合弱约束方向先验、有界导纳修正及下层跟踪能力，形成接触引导与受约束执行相协调的方法，提高典型接口在基座振动条件下的对准}精度与接触\rev{柔顺性}。

\subsection{年度研究计划及预期研究成果}

\subsubsection{年度研究计划}

\topic{2027年1月至6月}
围绕研究内容一\rev{，建立安装界面耦合作用的反向动力学描述，校核几何阻抗与动力学坐标映射；同步}开展研究内容三的\rev{安装端运动复现与多源观测，}确定激励\rev{和}负载工况，完善连接\rev{组件、测点}及执行时序。\rev{形成}动力学模型、\rev{界面基线}数据\rev{和实验评价}方案。

\topic{2027年7月至12月}
完成几何阻抗与动力学优化控制设计\rev{，分析建模偏差}及\rev{约束激活下的阻抗可实现性}，开展位姿保持、轨迹跟踪与阻抗响应\rev{实验}；启动接触观测\rev{误差分析和动态}接触场适用性\rev{检验}。形成论文初稿1篇\rev{及}年度报告，与重点实验室开展专题交流。

\topic{2028年1月至6月}
\rev{完成接触位姿局部可辨识性与误差传播分析，}研究\rev{弱约束方向先验及参考修正与驱动能力的协调}，\rev{将}PCF与RAC接入运动控制\rev{；}完成\rev{静止}基座四组对照\rev{及代表性}单向激励\rev{实验}，形成方法论文\rev{与}发明专利申请材料。

\topic{2028年7月至12月}
完成\rev{接触引导与受约束执行协调方法}验证，\rev{开展运动}与\rev{组装}集成\rev{实验，评价不同}基座\rev{激励和初始偏差}下的\rev{组装精度、接触载荷}及方法适用范围；整理程序、数据\rev{和}测试报告，完成论文发表或录用、专利申请\rev{、}成果交流\rev{与结题}。

\subsubsection{预期研究成果与学术交流}

对应三项研究内容，形成\rev{安装界面耦合动力学表征与受约束阻抗响应的分析结果，开发}包含\rev{反向}动力学、几何阻抗与动力学优化的机器人运动控制程序1套\rev{，以及}典型接口的接触配准与\rev{有界}参考修正模块1套\rev{；}建立\rev{安装端运动复现、同步观测与独立精度评价的}地面实验流程，形成同步测量数据集\rev{，}完成振动\rev{环境}下运动与典型接口组装\rev{的}地面原理验证，提交技术总结报告和\rev{实验}测试报告各1份。围绕新增研究发表或录用学术论文3篇\rev{，申请发明专利1项}，培养参与研究的研究生3名。

每年至少与重点实验室研究人员开展1次实质性学术交流，结合模型、实验和阶段成果组织研讨与学术报告；拟参加1次相关国内学术会议，并利用线上形式开展国际学术讨论。项目成果署名、资助标注和知识产权安排按重点实验室开放课题管理要求执行。

\section{研究基础与工作条件}

\subsection{工作基础}

团队围绕空间在轨操作，已形成柔性多体动力学、受约束机器人控制和接触对接的研究积累，可支撑本项目三项研究内容的衔接与实施。

\topic{（1）柔性多体动力学与模型简化基础}

申请人参与了柔性多体系统模态缩减研究\cite{sonneville2021}，团队进一步开展空间多柔体系统缩比等效建模\cite{qi2025}与几何精确梁动力学降阶研究，积累了柔性结构动力学分析和模型校核经验。上述研究积累可支撑本项目对机器人与柔性桁架相互作用的分析、安装端作用力与力矩的校核，以及地面试验条件与实际作业工况之间的对应分析。

\topic{（2）\rev{反向}动力学建模与运动控制基础}

团队针对七自由度作业机械臂，已实现以作业端为浮动基、通过实际安装端作用力与力矩描述外部作用的\rev{反向}动力学模型，并建立了几何阻抗与动力学优化相结合的运动控制程序。该程序利用关节状态以及运动捕捉系统与惯性测量单元（IMU）的测量数据生成运动参考，在关节运动、力矩及其变化约束下求取控制指令，已支持位姿保持、圆轨迹跟踪及固定基座轨迹跟踪试验，并具备数据记录与回放功能。现有安装端受力模型采用静力平衡近似，为本项目在模拟装置施加基座激励的条件下校核安装端动态作用力与力矩、研究末端运动与安装端力和力矩的关系提供了可扩展的实现基础。

\topic{（3）接触几何与柔顺对接基础}

团队已开展在轨组装柔顺控制与轨迹规划研究，并完成多模块对接力控制地面实验\cite{shi2024}。近期研究已形成PCF六维目标配准与RAC约束导纳捕获方法，建立了有界修正、零输入保持及参考端口离散广义功关系。论文主表选取的五轮真机试验中，完整方法在65次试验中58次满足两秒保持窗口内位置与姿态均方根误差不超过0.5~mm和0.5$^\circ$的判据，仅RAC为35/65；同组60段非零偏差配准回放的平均三维位置与姿态误差为2.085~mm和1.622$^\circ$。现有评分采用示教对齐基准，真机通过导纳与位置伺服执行。上述结果为本项目进一步开展动态基座下的力矩控制融合及独立测量验证提供了基础。

\subsection{工作条件}

\topic{（1）已具备的条件}

本项目依托现有基座振动模拟与机器人作业实验平台，通过基座运动模拟装置向作业机械臂安装端施加振动扰动。现有作业机械臂的开放力矩控制程序、运动捕捉系统与IMU的状态输入接口及同步记录接口，可支持运动控制集成与对照实验；已有模块对接平台、六维力/力矩测量及接触记录，可支撑接口标定和柔顺接触研究。作业机械臂的运动控制程序及接触记录回放流程为运动控制实现与接触对准方法验证提供了基础。

依托单位另建有约$4\,\mathrm m\times10\,\mathrm m$的在轨操作微重力模拟地面实验平台，可为相关地面验证提供配套条件。本项目悬挂装置的实验结果按照地面重力下的基座振动工况进行分析。

\topic{（2）尚需完善的条件及解决途径}

围绕基座振动下的闭环验证，拟完善模拟装置与作业机械臂之间的悬挂连接法兰、转接组件及可换配重，布置基座运动与安装端力和力矩测点，完成位姿测量系统、力和力矩传感器及传感器空间外参的标定，并实现模拟装置与作业机械臂的执行时序同步。结合现有对接平台完善典型接口夹具与独立位姿测量，使基座激励、末端运动和最终对接精度能够分别评价。连接组件、接口改装与专项测试按项目预算落实，计算与通用测量设备优先共享。

拟依托重点实验室在航空航天结构力学及控制方面的研究条件，围绕安装端作用力与力矩校核、运动与受力测量及控制实验开展联合讨论和测试，按阶段研究进度落实设备共享与试验安排。


\subsection{申请人简介}

\topic{（1）学历与研究工作简历}

单明贺，北京理工大学教授、博士生导师，国家级青年人才，主要从事空间在轨操作动力学与控制、柔性多体动力学和空间机器人技术研究。2007年9月至2011年7月在哈尔滨工业大学学习，获飞行器制造工程学士学位；2011年9月至2013年7月在哈尔滨工业大学学习，获航空宇航制造工程硕士学位；2013年11月至2018年6月在荷兰代尔夫特理工大学学习，获航天工程博士学位。2018年11月至2020年11月在美国马里兰大学帕克分校从事博士后研究，2021年1月起在北京理工大学工作。

申请人主持国家科技重大专项子课题、国家自然科学基金面上项目和青年项目等科研项目10余项，发表高水平论文30余篇，获授权发明专利7项。相关研究覆盖柔性结构动力学、空间机器人操作与地面验证，为本项目的耦合建模、控制方法研究和实验组织提供支撑。

\Needspace{7\baselineskip}
\topic{（2）近期与本项目有关的主要论著}

\begin{enumerate}[label=\arabic*.]
 \item \begin{minipage}[t]{\linewidth}祁伊凡, 单明贺, 田强. 空间多柔体系统缩比等效动力学建模方法研究[J]. 中国科学: 物理学 力学 天文学, 2025, 55(2): 224518. DOI: \doi{10.1360/SSPMA-2024-0302}.
 \end{minipage}
 \item \begin{minipage}[t]{\linewidth}史玲玲, 肖晓龙, 张晓峰, 范立佳, 单明贺, 田强. 空间机器人在轨组装多模块单元的对接力控制与地面实验[J]. 力学学报, 2024, 56(3): 800--816. DOI: \doi{10.6052/0459-1879-23-458}.
 \end{minipage}
 \item \begin{minipage}[t]{\linewidth}Yifan Qi, Shilei Han, Minghe Shan, Mingwu Li, Qiang Tian. Invariant-manifold-based model reduction for geometrically exact beam dynamics[J]. Computer Methods in Applied Mechanics and Engineering, 2026, 450: 118643. DOI: \doi{10.1016/j.cma.2025.118643}.
 \end{minipage}
\end{enumerate}

% 接触研究稿为匿名研究材料，作者名单未提供，已在“工作基础”说明。
% 不将匿名稿件冒列为申请人已发表论著，正式论著清单按实际发表状态更新。
\topic{（3）学术奖励}

\begin{enumerate}[label=\arabic*.]
\item 单明贺，国际空间研究委员会青年科学家最佳论文奖（COSPAR Outstanding Paper Award for Young Scientists），2022年。对应论文为：Minghe Shan, Jian Guo, Eberhard Gill. An analysis of the flexibility modeling of a net for space debris removal[J]. Advances in Space Research, 2020, 65(3): 1083--1094. DOI: \doi{10.1016/j.asr.2019.10.041}。
\item 张大羽，罗凯，赵将，韩石磊，朱佳龙，单明贺，王辉，刘凡，薛永刚，王立朋. 空间大型柔性可展结构高可靠展开与超低频模态辨识关键技术及应用. 中国振动工程学会科学技术奖二等奖，2023年度。
\end{enumerate}

\subsection{项目组主要成员简介}

项目组由单明贺、苏文康、黄朕、李融冰和任启明组成，围绕耦合建模、控制方法与地面验证拟作如下分工。

\textbf{单明贺：}项目负责人，研究基础及代表性成果见前节。负责总体研究方案与进度组织，统筹机器人与柔性桁架的耦合建模、闭环性能分析及与重点实验室的合作交流。

\textbf{苏文康：}已开展自由漂浮空间机器人任务优先协同操作与实时优化控制研究。拟承担几何阻抗与动力学优化控制设计、驱动约束协调及基座振动下的末端运动与安装端力和力矩分析。

\textbf{黄朕：}已开展大型空间望远镜多机器人组装序列规划研究，具备面向在轨组装任务组织操作流程的研究基础。拟承担典型作业过程与参考轨迹设计、实验工况组织及结果评价，服务于本项目单机器人运动与对接验证。

\textbf{李融冰：}拟承担作业机械臂的控制程序集成与实物实验，重点开展实时力矩接口实现、基座振动下运动控制验证和实验数据分析。

\textbf{\rev{任启明：}}\rev{拟承担接触配准与参考修正方法衔接、典型接口柔顺组装实验及误差分析。}

\Needspace{7\baselineskip}
\topic{近期与本项目有关的主要论著}

\begingroup
\small
\setstretch{1.12}
\begin{enumerate}[label=\arabic*.]
\item \begin{minipage}[t]{\linewidth}Wenkang Su, Lingling Shi, Andreas Müller, Minghe Shan. Cooperative manipulation control with task-prioritized real-time optimization for free-floating dual-arm space robots[J]. Aerospace Science and Technology, 2026, 171: 111584. DOI: \doi{10.1016/j.ast.2025.111584}.
\end{minipage}
\item \begin{minipage}[t]{\linewidth}Zhen Huang, Lingling Shi, Lijia Fan, Yan Yang, Minghe Shan. Multirobot Assembly Sequence Planning of a Large Space Telescope[J]. Journal of Spacecraft and Rockets, 2025, 62(3): 915--927. DOI: \doi{10.2514/1.A36129}.
\end{minipage}
\item \begin{minipage}[t]{\linewidth}Qiming Ren, Minghe Shan. A unified framework for compliant control and trajectory planning in robotic in-orbit assembly[J]. Acta Astronautica, 2026, 243: 32--45. DOI: \doi{10.1016/j.actaastro.2026.01.029}.
\end{minipage}
\end{enumerate}
\endgroup


\subsection{近五年承担科研项目情况}

本项目依托既有研究基础，重点研究柔性安装界面的动态作用、受约束阻抗响应及接触参考修正，形成相应模型、算法和典型接口实验结果。

\needinfo{近五年科研项目的正式名称、编号、经费来源、起止年月及承担角色，并按实际任务书说明与本项目的联系和任务区别。}

\clearpage
\section{经费预算}

% 重点课题20--30万元，本稿按30万元编制；数量和单价为需求测算，非供应商报价。
申请经费30.00万元，围绕模拟装置与作业机械臂的悬挂连接、接口改装、同步测量与控制接口、专项测试及成果产出安排。依托现有地面实验平台及计算设备开展研究，优先共享可用仪器设备。经费预算如下。

\begingroup
\small
\begin{longtable}{P{3.15cm}>{\centering\arraybackslash}p{1.35cm}P{9.6cm}}
\toprule
预算科目 & 金额（万元） & 计算根据及理由 \\
\midrule
\endfirsthead
\toprule
预算科目 & 金额（万元） & 计算根据及理由 \\
\midrule
\endhead
\midrule
\multicolumn{3}{r}{续下页}\\
\endfoot
\bottomrule
\endlastfoot
\textbf{合计} & \textbf{30.00} & 服务于耦合动力学与运动控制、接触组装和地面\rev{实验}验证三项研究内容。 \\
设备费 & 4.00 & 同步采集和实时控制接口改装所需设备及配件1套，按4.00万元测算，用于模拟装置与作业机械臂的状态、基座与末端运动及力和力矩测量数据的时序同步。 \\
材料费 & 6.00 & 悬挂连接法兰、转接组件及配套材料按3.60万元测算；典型接口夹具及可换配重材料1.40万元；传感安装、线缆等耗材1.00万元。 \\
测试化验加工费 & 9.00 & 悬挂连接及接口夹具加工装配3批，每批1.50万元，共4.50万元；基座与末端位姿测量、安装端力和力矩传感器标定3批，每批0.80万元，共2.40万元；外协动态测量与联合测试3批，每批0.70万元，共2.10万元。共享仪器免费使用部分不重复计费。 \\
差旅费 & 3.00 & 合作试验、专题交流及参加国内学术会议等约10人次，按每人次0.30万元测算。 \\
会议费 & 0.50 & 项目专题研讨及实验方案评议2次，按每次0.25万元测算。 \\
国际合作与交流费 & 0.00 & 以线上学术交流为主，不安排出国经费。 \\
出版/文献/信息传播/知识产权事务费 & 3.50 & 论文出版相关费用2篇，每篇按1.00万元测算，共2.00万元；发明专利申请及相关事务1.00万元；文献和信息资料0.50万元。 \\
专家咨询费 & 1.00 & 动力学建模及实验方案咨询约5人次，按每人次0.20万元测算；占总经费3.33\%。 \\
劳务费 & 3.00 & 参与建模、控制设计和实验工作的研究生2人，每人累计15个月，每人每月0.10万元；占总经费10.00\%。 \\
\end{longtable}
\endgroup

预算按研究需求测算，设备与测试项目结合现有条件核定，相关支出按实际发生及适用标准执行。专家咨询费和劳务费分别满足申请书模板规定的不超过总经费5\%和10\%的比例要求。

% ==================== 报告正文结束 ====================
\end{document}
